\chapter{2012年新课标卷（理）}

\section{单选题}

\begin{problem}
已知集合 \(A = \{1, 2, 3, 4, 5\}\)，\(B = \{(x, y) \mid x \in A, y \in A, x - y \in A\}\)，则 \(B\) 中所含元素的个数为（\quad）

\choices
    {\(3\)}
    {\(6\)}
    {\(8\)}
    {\(10\)}
\end{problem}

\begin{answer}
D.
\end{answer}

\begin{solution}
由 \(x-y\in A=\{1,2,3,4,5\}\) 可知 \(x>y\). 当 \(x-y=1\)，\(2\)，\(3\)，\(4\) 时，分别有 \(4\)，\(3\)，\(2\)，\(1\) 对 \(x\)，\(y\) 满足条件；差值不可能为 \(5\)，因为 \(x\)，\(y\in A\). 因此 \(B\) 中元素的个数为 \(4+3+2+1=10\)，故选 D.
\end{solution}

\begin{problem}
将 \(2\) 名教师，\(4\) 名学生分成 \(2\) 个小组，分别安排到甲，乙两地参加社会实践活动，每个小组由 \(1\) 名教师和 \(2\) 名学生组成，不同的安排方案共有（\quad）

\choices
    {\(12\text{ 种}\)}
    {\(10\text{ 种}\)}
    {\(9\text{ 种}\)}
    {\(8\text{ 种}\)}
\end{problem}

\begin{answer}
A.
\end{answer}

\begin{solution}
先把 \(2\) 名教师安排到甲、乙两地，每地各 \(1\) 名教师，有 \(2\) 种安排方法. 再从 \(4\) 名学生中选出 \(2\) 名安排到甲地，剩下的 \(2\) 名安排到乙地，有 \(\mathrm C_4^2=6\) 种方法. 因而不同的安排方案共有 \(2\times6=12\) 种，故选 A.
\end{solution}

\begin{problem}
下面是关于复数 \(z = \frac{2}{-1 + \i}\) 的四个命题：\(p_1: |z| = 2\)，\(p_2: z^2 = 2\i\)；\(p_3: z\) 的共轭复数为 \(1 + \i\)；\(p_4: z\) 的虚部为 \(-1\). 其中的真命题为（\quad）

\choices
    {\(p_{2}\)，\(p_{3}\)}
    {\(p_{1}\)，\(p_{2}\)}
    {\(p_{2}\)，\(p_{4}\)}
    {\(p_{3}\)，\(p_{4}\)}
\end{problem}

\begin{answer}
C.
\end{answer}

\begin{solution}
将分母有理化，得 \(z=\dfrac{2}{-1+\i}=\dfrac{2(-1-\i)}{(-1)^2+1^2}=-1-\i\). 因此 \(|z|=\sqrt{(-1)^2+(-1)^2}=\sqrt2\ne2\)，命题 \(p_1\) 为假；\(z^2=(-1-\i)^2=2\i\)，命题 \(p_2\) 为真；\(z\) 的共轭复数为 \(-1+\i\ne1+\i\)，命题 \(p_3\) 为假；\(z\) 的虚部为 \(-1\)，命题 \(p_4\) 为真. 所以真命题为 \(p_2\)，\(p_4\)，故选 C.
\end{solution}

\begin{problem}
设 \(F_{1}\)，\(F_{2}\) 是椭圆 \(E: \frac{x^{2}}{a^{2}} + \frac{y^{2}}{b^{2}} = 1\)（\(a > b > 0\)）的左，右焦点，\(P\) 为直线 \(x = \frac{3a}{2}\) 上一点，\(\triangle F_{2}PF_{1}\) 是底角为 \(30^{\circ}\) 的等腰三角形，则 \(E\) 的离心率为（\quad）

\choices
    {\( \frac{1}{2} \)}
    {\( \frac{2}{3} \)}
    {\( \frac{3}{4} \)}
    {\( \frac{4}{5} \)}
\end{problem}

\begin{answer}
C.
\end{answer}

\begin{solution}
设椭圆半焦距为 \(c\)，则 \(c^2=a^2-b^2\)，且焦点为 \(F_1=(-c,0)\)，\(F_2=(c,0)\). 由于 \(P\) 的横坐标为 \(\dfrac{3a}{2}\)，且 \(c<a<\dfrac{3a}{2}\)，点 \(P\) 在 \(F_2\) 的右侧，所以 \(\angle PF_2F_1>90^\circ\). 若等腰三角形的两腰为 \(PF_1\)，\(PF_2\)，则 \(P\) 在 \(F_1F_2\) 的垂直平分线上，与 \(x_P=\dfrac{3a}{2}\) 矛盾；若两腰为 \(PF_1\)，\(F_1F_2\)，则 \(\angle PF_2F_1=30^\circ\)，也与上式矛盾. 因此两腰为 \(PF_2\)，\(F_1F_2\)，从而 \(PF_2=F_1F_2=2c\)，且 \(\angle PF_1F_2=30^\circ\).

设 \(P=\left(\dfrac{3a}{2},t\right)\). 由 \(\angle PF_1F_2=30^\circ\)，有 \(|t|=\dfrac{1}{\sqrt3}\left(\dfrac{3a}{2}+c\right)\). 又由 \(PF_2=2c\)，有 \(\dfrac{1}{3}\left(\dfrac{3a}{2}+c\right)^2+\left(\dfrac{3a}{2}-c\right)^2=4c^2\). 令 \(r=\dfrac ca\)，化简得 \(8r^2+6r-9=0\). 因为 \(r>0\)，所以 \(r=\dfrac34\). 椭圆的离心率为 \(e=\dfrac ca=\dfrac34\)，故选 C.
\end{solution}

\begin{problem}
已知 \(\{a_{n}\}\) 为等比数列，\(a_{4} + a_{7} = 2\)，\(a_{5}a_{6} = -8\)，则 \(a_{1} + a_{10} =\)（\quad）

\choices
    {\(7\)}
    {\(5\)}
    {\(-5\)}
    {\(-7\)}
\end{problem}

\begin{answer}
D.
\end{answer}

\begin{solution}
设等比数列的公比为 \(q\). 因为 \(a_4a_7=a_1^2q^9=a_5a_6=-8\)，而 \(a_4+a_7=2\)，所以 \(a_4\)，\(a_7\) 是方程 \(t^2-2t-8=0\) 的两个根，即 \(\{a_4,a_7\}=\{4,-2\}\). 若 \(a_4=4\)，\(a_7=-2\)，则 \(q^3=\dfrac{a_7}{a_4}=-\dfrac12\)，从而 \(a_1=\dfrac{a_4}{q^3}=-8\)，\(a_{10}=a_7q^3=1\). 若 \(a_4=-2\)，\(a_7=4\)，则 \(q^3=-2\)，从而 \(a_1=1\)，\(a_{10}=-8\). 两种情形均有 \(a_1+a_{10}=-7\)，故选 D.
\end{solution}

\begin{problem}
如果执行下面的程序框图，输入正整数 \(N\)（\(N \geqslant 2\)）和实数 \(a_{1}\)，\(a_{2}\)，\(\cdots\)，\(a_{N}\)，输出 \(A\)，\(B\)，则（\quad）

\bitmapfigure[max width=0.46\linewidth]{img/2012/new_curriculum_science/q6_fig1.png}

\choices
    {\(A + B\) 为 \(a_{1}\)，\(a_{2}\)，\(\cdots\)，\(a_{N}\) 的和}
    {\(\frac{A + B}{2}\) 为 \(a_1\)，\(a_2\)，\(\dots\)，\(a_N\) 的算术平均数}
    {\(A\) 和 \(B\) 分别是 \(a_{1}\)，\(a_{2}\)，\(\dots\)，\(a_{N}\) 中最大的数和最小的数}
    {\(A\) 和 \(B\) 分别是 \(a_{1}\)，\(a_{2}\)，\(\dots\)，\(a_{N}\) 中最小的数和最大的数}
\end{problem}

\begin{answer}
C.
\end{answer}

\begin{solution}
初始时 \(k=1\)，且 \(A=B=a_1\). 每次循环先令 \(x=a_k\)，再比较 \(x\) 与当前的 \(A\)，\(B\)：若 \(x>A\)，就令 \(A=x\)；否则若 \(x<B\)，就令 \(B=x\). 因此处理完 \(a_1\)，\(a_2\)，\(\dots\)，\(a_k\) 后，\(A\) 始终是这些数中的最大值，\(B\) 始终是这些数中的最小值.

处理完当前的 \(a_k\) 后，程序判断 \(k\geqslant N\). 当 \(k<N\) 时，先令 \(k=k+1\)，再处理下一个数；当 \(k=N\) 时直接输出. 所以程序恰好比较完 \(a_1\)，\(a_2\)，\(\dots\)，\(a_N\)，输出的 \(A\)，\(B\) 分别为全部输入数中的最大值和最小值，故选 C.
\end{solution}

\begin{problem}
如图，网格上小正方形的边长为 \(1\)，粗线画出的是某几何体的三视图，则此几何体的体积为（\quad）

\bitmapfigure[max width=0.56\linewidth]{img/2012/new_curriculum_science/q7_fig1.png}

\choices
    {\(6\)}
    {\(9\)}
    {\(12\)}
    {\(18\)}
\end{problem}

\begin{answer}
B.
\end{answer}

\begin{solution}
由三视图可确定该几何体为三棱锥，俯视图给出其底面. 从网格数得底面是底边长为 \(6\)、高为 \(3\) 的等腰三角形，所以底面积为 \(S=\dfrac12\times6\times3=9\). 另外两个视图中竖直方向的最大投影长度均为 \(3\)，故棱锥的高为 \(h=3\). 因而几何体的体积为 \(V=\dfrac13Sh=\dfrac13\times9\times3=9\)，故选 B.
\end{solution}

\begin{problem}
等轴双曲线 \(C\) 的中心在原点，焦点在 \(x\) 轴上，\(C\) 与抛物线 \( y^{2} = 16x \) 的准线交于 \(A\)，\(B\) 两点，\( |AB| = 4\sqrt{3} \)，则 \(C\) 的实轴长为（\quad）

\choices
    {\(\sqrt{2}\)}
    {\(2\sqrt{2}\)}
    {\(4\)}
    {\(8\)}
\end{problem}

\begin{answer}
C.
\end{answer}

\begin{solution}
抛物线 \(y^2=16x\) 的准线为 \(x=-4\). 设等轴双曲线的半实轴为 \(r\)，则其方程可写为 \(\dfrac{x^2}{r^2}-\dfrac{y^2}{r^2}=1\). 双曲线与准线的交点 \(A\)，\(B\) 的横坐标均为 \(-4\)，代入得 \(y^2=16-r^2\). 因而 \(|AB|=2\sqrt{16-r^2}=4\sqrt3\)，解得 \(r^2=4\)，\(r=2\). 所以双曲线的实轴长为 \(2r=4\)，故选 C.
\end{solution}

\begin{problem}
已知 \(\omega > 0\)，函数 \(f(x) = \sin \left(\omega x + \frac{\pi}{4}\right)\) 在 \(\left(\frac{\pi}{2}, \pi\right)\) 单调递减，则 \(\omega\) 的取值范围是（\quad）

\choices
    {\(\left[\frac{1}{2}, \frac{5}{4}\right]\)}
    {\(\left[\frac{1}{2}, \frac{3}{4}\right]\)}
    {\(\left(0, \frac{1}{2}\right]\)}
    {\((0, 2]\)}
\end{problem}

\begin{answer}
A.
\end{answer}

\begin{solution}
函数的导数为 \(f'(x)=\omega\cos\left(\omega x+\dfrac\pi4\right)\). 因为 \(\omega>0\)，要使函数在整个 \(\left(\dfrac\pi2,\pi\right)\) 上单调递减，需有 \(\cos\left(\omega x+\dfrac\pi4\right)\leqslant0\) 对该区间内所有 \(x\) 成立. 令相位 \(t=\omega x+\dfrac\pi4\)，则相位区间为 \(I=\left(\dfrac{\omega\pi}{2}+\dfrac\pi4,\omega\pi+\dfrac\pi4\right)\). 因而存在整数 \(k\)，使 \(I\subseteq\left[\dfrac\pi2+2k\pi,\dfrac{3\pi}{2}+2k\pi\right]\). 由于 \(I\) 的左端大于 \(\dfrac\pi4\)，只能有 \(k\geqslant0\). 于是

\(\frac{\omega\pi}{2}+\frac\pi4\geqslant\frac\pi2+2k\pi\)，\(\omega\pi+\frac\pi4\leqslant\frac{3\pi}{2}+2k\pi\)

即 \(\omega\geqslant\dfrac12+4k\)，\(\omega\leqslant\dfrac54+2k\). 当 \(k\geqslant1\) 时，\(\dfrac12+4k>\dfrac54+2k\)，不可能；所以 \(k=0\). 因此 \(\dfrac12\leqslant\omega\leqslant\dfrac54\)，即 \(\omega\in\left[\dfrac12,\dfrac54\right]\)，故选 A.
\end{solution}

\begin{problem}
已知函数 \( f(x) = \frac{1}{\ln (x + 1) - x} \)，则 \( y = f(x) \) 的图象大致为（\quad）

\choices
    {\choicebitmap[width=0.15\paperwidth]{img/2012/new_curriculum_science/q10_optA.png}}
    {\choicebitmap[width=0.15\paperwidth]{img/2012/new_curriculum_science/q10_optB.png}}
    {\choicebitmap[width=0.15\paperwidth]{img/2012/new_curriculum_science/q10_optC.png}}
    {\choicebitmap[width=0.15\paperwidth]{img/2012/new_curriculum_science/q10_optD.png}}
\end{problem}

\begin{answer}
B.
\end{answer}

\begin{solution}
令 \(g(x)=\ln(x+1)-x\)，则 \(f\) 的定义域为 \(x>-1\) 且 \(x\ne0\). 由 \(g'(x)=\dfrac1{x+1}-1=-\dfrac{x}{x+1}\)，可知 \(g\) 在 \((-1,0)\) 上递增，在 \((0,+\infty)\) 上递减，且 \(g(0)=0\). 因为不等式 \(\ln(1+x)<x\) 在 \(x\ne0\) 且 \(x>-1\) 时成立，所以在定义域内始终有 \(g(x)<0\). 当 \(x\to-1^+\) 或 \(x\to+\infty\) 时，\(f(x)\to0^-\)；当 \(x\to0^-\) 或 \(x\to0^+\) 时，\(f(x)\to-\infty\). 因此图象有两支，均在 \(x\) 轴下方，\(x=0\) 是竖直渐近线，且两端以 \(y=0\) 为水平渐近线，符合选项 B，故选 B.
\end{solution}

\begin{problem}
已知三棱锥 \(S - ABC\) 的所有顶点都在球 \(O\) 的球面上，\(\triangle ABC\) 是边长为 \(1\) 的正三角形，\(SC\) 为球 \(O\) 的直径，且 \(SC = 2\)，则此棱锥的体积为（\quad）

\choices
    {\(\frac{\sqrt{2}}{6}\)}
    {\( \frac{\sqrt{3}}{6} \)}
    {\( \frac{\sqrt{2}}{3} \)}
    {\( \frac{\sqrt{2}}{2} \)}
\end{problem}

\begin{answer}
A.
\end{answer}

\begin{solution}
设球的半径为 \(R=1\)，球心为 \(O\)，则 \(SC=2\) 且 \(O\) 是 \(SC\) 的中点. 设平面 \(ABC\) 与球面的截面圆圆心为 \(O_1\). 因为 \(ABC\) 是边长为 \(1\) 的正三角形，其外接圆半径为 \(O_1A=\dfrac{1}{\sqrt3}=\dfrac{\sqrt3}{3}\). 由直角三角形 \(OO_1A\)，得 \(OO_1=\sqrt{OA^2-O_1A^2}=\sqrt{1-\dfrac13}=\sqrt{\dfrac23}=\dfrac{\sqrt6}{3}\). 由于 \(S\)，\(O\)，\(C\) 共线且 \(O\) 是 \(SC\) 的中点，点 \(S\) 到平面 \(ABC\) 的距离为 \(2OO_1=\dfrac{2\sqrt6}{3}\). 又 \(S_{\triangle ABC}=\dfrac{\sqrt3}{4}\)，所以 \(V_{S-ABC}=\dfrac13S_{\triangle ABC}\cdot\dfrac{2\sqrt6}{3}=\dfrac13\times\dfrac{\sqrt3}{4}\times\dfrac{2\sqrt6}{3}=\dfrac{\sqrt2}{6}\). 故选 A.
\end{solution}

\begin{problem}
设点 \(P\) 在曲线 \(y = \frac{1}{2} \e^{x}\) 上，点 \(Q\) 在曲线 \(y = \ln (2x)\) 上，则 \(|PQ|\) 的最小值为（\quad）

\choices
    {\(1 - \ln 2\)}
    {\(\sqrt{2}\,(1 - \ln 2)\)}
    {\(1 + \ln 2\)}
    {\(\sqrt{2}\,(1 + \ln 2)\)}
\end{problem}

\begin{answer}
B.
\end{answer}

\begin{solution}
曲线 \(y=\dfrac12\e^x\) 可改写为 \(x=\ln(2y)\)，所以它与曲线 \(y=\ln(2x)\) 关于直线 \(y=x\) 对称. 对第一条曲线上的点 \(P=(x,y)\)，有 \(y-x=\dfrac12\e^x-x\). 令 \(h(x)=\dfrac12\e^x-x\)，则 \(h'(x)=\dfrac12\e^x-1\)，故 \(h\) 在 \(x=\ln2\) 处取得最小值 \(h(\ln2)=1-\ln2\). 因而第一条曲线上的点都满足 \(y-x\geqslant1-\ln2\). 同理，第二条曲线上的点 \(Q=(u,v)\) 都满足 \(u-v\geqslant1-\ln2\). 点 \(P\) 与点 \(Q\) 到直线 \(y=x\) 的有向距离之差至少为 \(\dfrac{1-\ln2}{\sqrt2}+\dfrac{1-\ln2}{\sqrt2}=\sqrt2(1-\ln2)\)，所以 \(|PQ|\geqslant\sqrt2(1-\ln2)\). 当 \(P=(\ln2,1)\)，\(Q=(1,\ln2)\) 时，两点分别在两条曲线上，且连线垂直于 \(y=x\)，达到上述下界. 因此最小值为 \(\sqrt2(1-\ln2)\)，故选 B.
\end{solution}

\section{填空题}

\begin{problem}
已知向量 \(\bs{a}\)，\(\bs{b}\) 夹角为 \(45^{\circ}\)，且 \(|\bs{a}| = 1\)，\(|2\bs{a} - \bs{b}| = \sqrt{10}\)，则 \(|\bs{b}| =\fillinblank{}\).
\end{problem}

\begin{answer}
\(3\sqrt{2}\).
\end{answer}

\begin{solution}
由 \(|2\bs{a}-\bs{b}|^2=4|\bs{a}|^2+|\bs{b}|^2-4\bs{a}\cdot\bs{b}\)，且 \(\bs{a}\cdot\bs{b}=|\bs{a}||\bs{b}|\cos45^\circ=\dfrac{|\bs{b}|}{\sqrt2}\)，得 \(10=4+|\bs{b}|^2-2\sqrt2|\bs{b}|\). 令 \(t=|\bs{b}|>0\)，则 \(t^2-2\sqrt2t-6=0\)，即 \((t-3\sqrt2)(t+\sqrt2)=0\). 因为 \(t>0\)，所以 \(|\bs{b}|=3\sqrt2\).
\end{solution}

\begin{problem}
设 \(x\)，\(y\) 满足约束条件 \(\begin{cases}
x - y \geqslant -1,\\
x + y \leqslant 3,\\
x \geqslant 0,\\
y \geqslant 0,
\end{cases}\) 则 \(z = x - 2y\) 的取值范围为\fillinblank{}.
\end{problem}

\begin{answer}
\(\left[-3,3\right]\).
\end{answer}

\begin{solution}
可行域由 \(x\geqslant0\)，\(y\geqslant0\)，\(y\leqslant x+1\)，\(y\leqslant3-x\) 围成. 其顶点为 \((0,0)\)，\((0,1)\)，\((1,2)\)，\((3,0)\). 目标函数 \(z=x-2y\) 在这些顶点处的值依次为 \(0\)，\(-2\)，\(-3\)，\(3\). 由于线性函数在有界多边形上的最大值和最小值均在顶点取得，所以 \(z\) 的取值范围为 \([-3,3]\).
\end{solution}

\begin{problem}
某一部件由三个电子元件按下图方式连接而成，元件 \(1\) 或元件 \(2\) 正常工作，且元件 \(3\) 正常工作，则部件正常工作. 设三个电子元件的使用寿命（单位：小时）均服从正态分布 \(N(1000, 50^{2})\)，且各个元件能否正常工作相互独立，那么该部件的使用寿命超过 \(1000\) 小时的概率为\fillinblank{}.

\bitmapfigure[max width=0.52\linewidth]{img/2012/new_curriculum_science/q15_fig1.png}
\end{problem}

\begin{answer}
\(\frac{3}{8}\).
\end{answer}

\begin{solution}
设三个元件在时刻 \(1000\) 小时后仍正常工作的事件分别为 \(E_1\)，\(E_2\)，\(E_3\). 由于每个元件寿命服从以 \(1000\) 为均值的正态分布，分布关于 \(1000\) 对称，所以 \(P(E_i)=\dfrac12\). 电路中元件 \(1\)，\(2\) 并联，元件 \(3\) 与其串联，故部件正常工作的事件为 \((E_1\cup E_2)\cap E_3\). 由独立性，\(P(E_1\cup E_2)=1-P(\overline E_1\cap\overline E_2)=1-\dfrac14=\dfrac34\)，从而所求概率为 \(P((E_1\cup E_2)\cap E_3)=\dfrac34\times\dfrac12=\dfrac38\).
\end{solution}

\begin{problem}
数列 \(\{a_{n}\}\) 满足 \(a_{n + 1} + (-1)^n a_n = 2n - 1\)，则 \(\{a_n\}\) 的前 \(60\) 项和为\fillinblank{}.
\end{problem}

\begin{answer}
\(1830\).
\end{answer}

\begin{solution}
令 \(S_{60}=\sum_{n=1}^{60}a_n\). 对 \(n=2k-1\) 和 \(n=2k\)，递推式分别给出 \(a_{2k}-a_{2k-1}=4k-3\) 与 \(a_{2k+1}+a_{2k}=4k-1\). 消去 \(a_{2k}\)，得到 \(a_{2k+1}=2-a_{2k-1}\). 因此 \(a_{2k-1}=a_1\)（\(k\) 为奇数），\(a_{2k-1}=2-a_1\)（\(k\) 为偶数）. 对 \(j=1\)，\(\dots\)，\(15\)，有
\(a_{4j-3}+a_{4j-2}=2a_1+8j-7\)，\(a_{4j-1}+a_{4j}=1-2a_1+8j\). 所以 \(S_{60}=\sum_{j=1}^{15}(16j-6)=16\cdot\dfrac{15\cdot16}{2}-6\cdot15=1830\).
\end{solution}

\section{解答题}

\begin{problem}
已知 \(a\)，\(b\)，\(c\) 分别为 \(\triangle ABC\) 三个内角 \(A\)，\(B\)，\(C\) 的对边，\(a \cos C + \sqrt{3} a \sin C - b - c = 0\).

\begin{enumerate}
\item 求 \(A\)；
\item 若 \(a = 2\)，\(\triangle ABC\) 的面积为 \(\sqrt{3}\)，求 \(b\)，\(c\).
\end{enumerate}
\end{problem}

\begin{solution}
\begin{enumerate}
\item 由 \(B=\pi-A-C\)，正弦定理给出 \(b/a=\sin B/\sin A\)，\(c/a=\sin C/\sin A\). 将题设等式除以 \(a\)，并乘以 \(\sin A\)，得 \(\sin A\cos C+\sqrt3\sin A\sin C-\sin B-\sin C=0\). 又 \(\sin B=\sin(A+C)=\sin A\cos C+\cos A\sin C\)，所以 \(\sin C(\sqrt3\sin A-\cos A-1)=0\). 因为 \(0<C<\pi\)，有 \(\sin C>0\)，故 \(\sqrt3\sin A-\cos A=1\). 即 \(2\sin\left(A-\dfrac\pi6\right)=1\). 由于 \(0<A<\pi\)，只有 \(A=\dfrac\pi3=60^\circ\) 满足条件.
\item 三角形面积为 \(\dfrac12bc\sin A=\sqrt3\)，代入 \(A=60^\circ\) 得 \(bc=4\). 由余弦定理 \(a^2=b^2+c^2-2bc\cos A\)，代入 \(a=2\)，\(A=60^\circ\) 和 \(bc=4\)，得 \(b^2+c^2=8\). 因而 \((b+c)^2=b^2+c^2+2bc=16\)，且 \(b+c=4\). 所以 \(b\)，\(c\) 是方程 \(t^2-4t+4=0\) 的两个正根，故 \(b=c=2\).
\end{enumerate}
\end{solution}

\begin{problem}
某花店每天以每枝 \(5\) 元的价格从农场购进若干枝玫瑰花，然后以每枝 \(10\) 元的价格出售，如果当天卖不完，剩下的玫瑰花作垃圾处理.

\begin{enumerate}
\item 若花店一天购进 \(16\) 枝玫瑰花，求当天的利润 \(y\)（单位：元）关于当天需求量 \(n\)（单位：枝，\(n \in \N\)）的函数解析式；
\item 花店记录了 \(100\) 天玫瑰花的日需求量（单位：枝），整理得下表：

\begin{center}
\renewcommand{\arraystretch}{1.2}
\begin{tabular}{|c|c|c|c|c|c|c|c|}
\hline
日需求量 \(n\) & \(14\) & \(15\) & \(16\) & \(17\) & \(18\) & \(19\) & \(20\) \\
\hline
频数 & \(10\) & \(20\) & \(16\) & \(16\) & \(15\) & \(13\) & \(10\) \\
\hline
\end{tabular}
\end{center}

以 \(100\) 天记录的各需求量的频率作为各需求量发生的概率.

\begin{enumerate}
\item 若花店一天购进 \(16\) 枝玫瑰花，\(X\) 表示当天的利润（单位：元），求 \(X\) 的分布列，数学期望及方差；
\item 若花店计划一天购进 \(16\) 枝或 \(17\) 枝玫瑰花，你认为应购进 \(16\) 枝还是 \(17\) 枝？请说明理由.
\end{enumerate}
\end{enumerate}
\end{problem}

\begin{solution}
\begin{enumerate}
\item 一天购进 \(16\) 枝时，购进成本为 \(80\) 元，实际售出 \(\min(n,16)\) 枝，故利润为 \(10\min(n,16)-80\). 因而
\[
y=
\begin{cases}
10n-80,&n\leqslant16,\\
80,&n>16.
\end{cases}
\]
\item
\begin{enumerate}
\item 当购进 \(16\) 枝时，需求量 \(n=14\)，\(15\) 时利润分别为 \(60\)，\(70\) 元，需求量 \(n=16\)，\(17\)，\(18\)，\(19\)，\(20\) 时利润均为 \(80\) 元. 由频数除以 \(100\)，得 \(P(X=60)=0.1\)，\(P(X=70)=0.2\)，\(P(X=80)=0.7\). 因而 \(E(X)=60\times0.1+70\times0.2+80\times0.7=76\)，\(E(X^2)=60^2\times0.1+70^2\times0.2+80^2\times0.7=5820\)，所以 \(D(X)=E(X^2)-[E(X)]^2=5820-76^2=44\).
\item 当购进 \(17\) 枝时，需求量 \(n=14\)，\(15\)，\(16\) 时利润分别为 \(55\)，\(65\)，\(75\) 元，需求量 \(n\geqslant17\) 时利润均为 \(85\) 元. 所以 \(E(Y)=55\times0.1+65\times0.2+75\times0.16+85\times0.54=76.4\). 因为 \(76.4>76\)，按数学期望最大原则，应购进 \(17\) 枝.
\end{enumerate}
\end{enumerate}
\end{solution}

\begin{problem}
如图，直三棱柱 \(ABC - A_{1}B_{1}C_{1}\) 中，\(AC = BC = \frac{1}{2} AA_{1}\)，\(D\) 是棱 \(AA_{1}\) 的中点，\(DC_{1} \perp BD\).

\begin{enumerate}
\item 证明：\(DC_{1} \perp BC\)；
\item 求二面角 \(A_{1} - BD - C_{1}\) 的大小.
\end{enumerate}
\bitmapfigure[max width=0.45\linewidth]{img/2012/new_curriculum_science/q19_fig1.png}
\end{problem}

\begin{solution}
\begin{enumerate}
\item 设 \(AA_1=h\)，令 \(r=\dfrac h2\). 取 \(A\) 为原点，使 \(AA_1\) 为 \(z\) 轴正方向，并设 \(\overrightarrow{AC}=\bs{u}\)，\(\overrightarrow{AB}=\bs{v}\). 则 \(D\) 的位置向量为 \(r\bs{k}\)，\(C_1\) 的位置向量为 \(\bs{u}+2r\bs{k}\)，\(B\) 的位置向量为 \(\bs{v}\). 由 \(DC_1\perp BD\)，有 \((\bs{u}+r\bs{k})\cdot(\bs{v}-r\bs{k})=0\)，即 \(\bs{u}\cdot\bs{v}=r^2\). 又 \(|\bs{u}|=AC=r\)，\(|\bs{v}-\bs{u}|=BC=r\)，所以 \(\bs{u}\cdot(\bs{v}-\bs{u})=0\)，从而 \(AC\perp BC\). 由于 \(\overrightarrow{BC}=\bs{u}-\bs{v}\)，\(\overrightarrow{DC_1}=\bs{u}+r\bs{k}\)，且 \(\bs{u}\)，\(\bs{v}\) 均垂直于 \(\bs{k}\)，有 \(\overrightarrow{BC}\cdot\overrightarrow{DC_1}=|\bs{u}|^2-\bs{u}\cdot\bs{v}=r^2-r^2=0\)，即 \(DC_1\perp BC\).
\item 在上述坐标系中可进一步取 \(A=(0,0,0)\)，\(C=(r,0,0)\)，\(B=(r,r,0)\)，\(A_1=(0,0,2r)\)，\(D=(0,0,r)\)，\(C_1=(r,0,2r)\). 平面 \(A_1BD\) 的一个法向量可取 \(\bs{n}_1=\overrightarrow{DB}\times\overrightarrow{DA_1}=r^2(1,-1,0)\)，平面 \(C_1BD\) 的一个法向量可取 \(\bs{n}_2=\overrightarrow{DB}\times\overrightarrow{DC_1}=r^2(1,-2,-1)\). 因而
\(\cos\theta=\dfrac{\bs{n}_1\cdot\bs{n}_2}{|\bs{n}_1||\bs{n}_2|}=\dfrac{3}{\sqrt2\sqrt6}=\dfrac{\sqrt3}{2}\). 二面角取锐角，故 \(\theta=30^\circ\).
\end{enumerate}
\end{solution}

\begin{problem}
设抛物线 \(C: x^{2} = 2py\)（\(p > 0\)）的焦点为 \(F\)，准线为 \(l\)，\(A\) 为 \(C\) 上一点，已知以 \(F\) 为圆心，\(FA\) 为半径的圆 \(F\) 交 \(l\) 于 \(B\)，\(D\) 两点.

\begin{enumerate}
\item 若 \(\angle BFD = 90^{\circ}\)，\(\triangle ABD\) 的面积为 \(4\sqrt{2}\)，求 \(p\) 的值及圆 \(F\) 的方程；
\item 若 \(A\)，\(B\)，\(F\) 三点在同一直线 \(m\) 上，直线 \(n\) 与 \(m\) 平行，且 \(n\) 与 \(C\) 只有一个公共点，求坐标原点到 \(m\)，\(n\) 距离的比值.
\end{enumerate}
\end{problem}

\begin{solution}
\begin{enumerate}
\item 抛物线 \(x^2=2py\) 的焦点为 \(F=(0,\dfrac p2)\)，准线为 \(l:y=-\dfrac p2\). 设 \(B=(q,-\dfrac p2)\)，\(D=(-q,-\dfrac p2)\). 因为 \(B\)，\(D\) 在以 \(F\) 为圆心的同一圆上，且 \(\angle BFD=90^\circ\)，有 \((q,-p)\cdot(-q,-p)=0\)，所以 \(q=p\). 因而 \(BD=2p\)，且圆的半径 \(FB=\sqrt2p\)，其方程为 \(x^2+\left(y-\dfrac p2\right)^2=2p^2\). 设 \(A=(x,y)\)，由 \(A\) 在抛物线上得 \(x^2=2py\). 又 \(A\) 在该圆上，故 \(2py+\left(y-\dfrac p2\right)^2=2p^2\)，即 \(y^2+py-\dfrac74p^2=0\). 因为 \(y>0\)，所以 \(y=p\left(\sqrt2-\dfrac12\right)\). 三角形 \(ABD\) 以 \(BD\) 为底，点 \(A\) 到 \(l\) 的高为 \(y+\dfrac p2=p\sqrt2\)，故其面积为 \(p^2\sqrt2=4\sqrt2\). 所以 \(p=2\)，圆 \(F\) 的方程为 \(x^2+(y-1)^2=8\).
\item 由 \(A\)，\(B\)，\(F\) 三点共线且 \(FA=FB\)，点 \(B\) 是点 \(A\) 关于 \(F\) 的中心对称点. 设 \(A=(x,y)\)，则 \(B=2F-A=(-x,p-y)\). 因为 \(B\in l\)，有 \(p-y=-\dfrac p2\)，故 \(y=\dfrac{3p}{2}\). 再由 \(A\in C\)，得 \(x^2=3p^2\)，所以直线 \(m\) 的斜率为 \(\dfrac1{\sqrt3}\) 或 \(-\dfrac1{\sqrt3}\)，且其到原点的距离为 \(\dfrac{\sqrt3p}{4}\). 与 \(m\) 平行的直线 \(n\) 可写为 \(y=\dfrac{x}{\sqrt3}+b\) 或 \(y=-\dfrac{x}{\sqrt3}+b\). 对第一种斜率，代入 \(x^2=2py\) 得

\[
x^2-\frac{2p}{\sqrt3}x-2pb=0
\]

其判别式为 \(\dfrac{4p^2}{3}+8pb\)，故相切时 \(b=-\dfrac p6\). 对第二种斜率同理得到

\[
x^2+\frac{2p}{\sqrt3}x-2pb=0
\]

判别式为 \(\dfrac{4p^2}{3}+8pb\)，同样有 \(b=-\dfrac p6\). 因而原点到 \(n\) 的距离为 \(\dfrac{\sqrt3p}{12}\)，所以所求距离之比为 \(3:1\).
\end{enumerate}
\end{solution}

\begin{problem}
已知函数 \(f(x)\) 满足 \(f(x) = f'(1)\e^{x-1} - f(0)x + \frac{1}{2}x^2\).

\begin{enumerate}
\item 求 \(f(x)\) 的解析式及单调区间；
\item 若 \(f(x) \geqslant \frac{1}{2}x^2 + ax + b\)，求 \((a + 1)b\) 的最大值.
\end{enumerate}
\end{problem}

\begin{solution}
\begin{enumerate}
\item 记 \(A=f'(1)\)，\(B=f(0)\). 在原式中令 \(x=0\)，得 \(B=A\e^{-1}\). 对原式求导，得 \(f'(x)=A\e^{x-1}-B+x\). 再令 \(x=1\)，结合 \(A=f'(1)\)，得 \(A=A-B+1\)，所以 \(B=1\)，\(A=\e\). 因而
\(f(x)=\e^x-x+\frac{x^2}{2}\)，\(f'(x)=\e^x+x-1\).
又 \(f''(x)=\e^x+1>0\)，且 \(f'(0)=0\)，所以 \(f'(x)<0\)（\(x<0\)），\(f'(x)>0\)（\(x>0\)）. 因此 \(f(x)\) 在 \((-\infty,0)\) 上单调递减，在 \((0,+\infty)\) 上单调递增.
\item 原不等式等价于
\(\e^x-(a+1)x-b\geqslant0\)，\((x\in\mathbb R)\).
令 \(t=a+1\). 若 \(t<0\)，令 \(x\to-\infty\)，左边趋于 \(-\infty\)，不可能对所有 \(x\) 成立. 若 \(t=0\)，需 \(b\leqslant0\)，于是 \(tb=0\). 若 \(t>0\)，函数 \(\e^x-tx\) 在 \(x=\ln t\) 处取得最小值 \(t(1-\ln t)\)，故
\(b\leqslant t(1-\ln t)\)，\(tb\leqslant t^2(1-\ln t)\).
令 \(H(t)=t^2(1-\ln t)\)，则 \(H'(t)=t(1-2\ln t)\). 因而 \(H\) 在 \(t=\sqrt{\e}\) 处取得最大值
\[
H(\sqrt{\e})=\frac{\e}{2}
\]
取 \(t=\sqrt{\e}\)，\(b=t(1-\ln t)=\frac{\sqrt{\e}}2\) 可达到该值，所以 \((a+1)b\) 的最大值为 \(\frac{\e}{2}\).
\end{enumerate}
\end{solution}

\section{选考题：解答题}

\begin{problem}
如图，\(D\)，\(E\) 分别为 \(\triangle ABC\) 边 \(AB\)，\(AC\) 的中点，直线 \(DE\) 交 \(\triangle ABC\) 的外接圆于 \(F\)，\(G\) 两点. 若 \(CF \parallel AB\)，证明：

\begin{enumerate}
\item \(CD = BC\)；
\item \(\triangle BCD \sim \triangle GBD\).
\end{enumerate}
\bitmapfigure[max width=0.36\linewidth]{img/2012/new_curriculum_science/q22_fig1.png}
\end{problem}

\begin{solution}
\begin{enumerate}
\item 由 \(D\)，\(E\) 为中点，得 \(DE\parallel BC\). 取坐标系 \(B=(0,0)\)，\(C=(2,0)\)，\(A=(u,v)\)，其中 \(v\ne0\). 则
\(D=\left(\frac u2,\frac v2\right)\)，\(F=\left(2+\frac u2,\frac v2\right)\)
后式由 \(F\in DE\) 且 \(CF\parallel AB\) 得到. 设外接圆方程为 \(x^2+y^2-2x+\lambda y=0\). 将 \(A\)、\(F\) 代入并消去 \(\lambda\)，得 \(u^2+v^2=8u\). 因而
\[
CD^2=\left(\frac u2-2\right)^2+\frac{v^2}{4}
=\frac{u^2+v^2-8u+16}{4}=4=BC^2
\]
所以 \(CD=BC\).
\item 直线 \(DE\) 为 \(y=v/2\). 圆与该直线的两个交点横坐标之和为 \(2\)，一个交点为 \(F\) 的横坐标 \(2+u/2\)，故另一个交点
\[
G=\left(-\frac u2,\frac v2\right)
\]
由 \(u^2+v^2=8u\)，可得
\(BD=GB=\sqrt{2u}\)，\(GD=u\)，\(BC=CD=2\).
于是 \(BC:GB=CD:BD=BD:GD\)，三边对应成比例，故 \(\triangle BCD\sim\triangle GBD\).
\end{enumerate}
\end{solution}

\begin{problem}
已知曲线 \(C_1\) 的参数方程是 \(\begin{cases}
x = 2\cos \varphi,\\
y = 3\sin \varphi,
\end{cases}\) （\(\varphi\) 是参数）以坐标原点为极点，\(x\) 轴的非负半轴为极轴建立极坐标系，曲线 \(C_2\) 的极坐标方程是 \(\rho = 2\)，正方形 \(ABCD\) 的顶点都在 \(C_2\) 上，且 \(A\)，\(B\)，\(C\)，\(D\) 依逆时针次序排列，点 \(A\) 的极坐标为 \(\left(2, \frac{\pi}{3}\right)\).

\begin{enumerate}
\item 求点 \(A\)，\(B\)，\(C\)，\(D\) 的直角坐标；
\item 设 \(P\) 为 \(C_1\) 上任意一点，求 \(|PA|^2 + |PB|^2 + |PC|^2 + |PD|^2\) 的取值范围.
\end{enumerate}
\end{problem}

\begin{solution}
\begin{enumerate}
\item 极坐标到直角坐标的转换公式为 \(x=\rho\cos\theta\)，\(y=\rho\sin\theta\). 因而
\[
A=\left(1,\sqrt3\right)
\]
正方形内角所对圆心角均为 \(\pi/2\)，依逆时针旋转 \(\pi/2\) 得
\(B=(-\sqrt3,1)\)，\(C=(-1,-\sqrt3)\)，\(D=(\sqrt3,-1)\).
\item 设 \(P=(x_P,y_P)=(2\cos\varphi,3\sin\varphi)\). 正方形中心为原点，故
\(A+B+C+D=(0,0)\)，\(|A|^2=|B|^2=|C|^2=|D|^2=4\).
展开平方和得
\[
\begin{aligned}
\sum_{V\in\{A,B,C,D\}}|P-V|^2
&=4(x_P^2+y_P^2)-2P\cdot(A+B+C+D)\\
&\quad+\sum_{V\in\{A,B,C,D\}}|V|^2\\
&=4(x_P^2+y_P^2)+16\\
&=16\cos^2\varphi+36\sin^2\varphi+16\\
&=32+20\sin^2\varphi.
\end{aligned}
\]
因为 \(\sin^2\varphi\in[0,1]\)，取值范围为 \([32,52]\).
\end{enumerate}
\end{solution}

\begin{problem}
已知函数 \( f(x) = |x + a| + |x - 2| \).

\begin{enumerate}
\item 当 \(a = -3\) 时，求不等式 \(f(x) \geqslant 3\) 的解集；
\item 若 \( f(x) \leqslant |x - 4| \) 的解集包含 \([1, 2]\)，求 \( a \) 的取值范围.
\end{enumerate}
\end{problem}

\begin{solution}
\begin{enumerate}
\item 当 \(a=-3\) 时，\(f(x)=|x-3|+|x-2|\). 分段讨论：若 \(x\leqslant2\)，则 \(f(x)=5-2x\)，不等式给出 \(x\leqslant1\)；若 \(2\leqslant x\leqslant3\)，则 \(f(x)=1<3\)，无解；若 \(x\geqslant3\)，则 \(f(x)=2x-5\)，不等式给出 \(x\geqslant4\). 所以解集为 \((-\infty,1]\cup[4,+\infty)\).
\item 对任意 \(x\in[1,2]\)，有 \(|x-2|=2-x\)，\(|x-4|=4-x\). 所以题设包含条件等价于
\[
|x+a|+2-x\leqslant4-x\quad\text{对所有 }x\in[1,2]
\]
即 \(|x+a|\leqslant2\) 对所有 \(x\in[1,2]\) 成立. 因为 \(x+a\) 在该区间的取值区间为 \([a+1,a+2]\)，这要求 \([a+1,a+2]\subseteq[-2,2]\). 因而 \(a\geqslant-3\)，\(a\leqslant0\)，即 \(a\in[-3,0]\). 反之该范围内的 \(a\) 确实满足条件.
\end{enumerate}
\end{solution}
